\chapter{Batch Clearing Algorithm}
\label{app:algorithm}

The clearing routine implements \eqref{eq:demandsupply} and
\eqref{eq:clearing} by a search over the admissible price grid, followed
by the tie-break cascade and the rationing step of
\Cref{sec:rationing}.

\begin{lstlisting}[language=,caption={Batch clearing, single asset
pair.},label={lst:clearing}]
INPUT : buys B = {(pi_i, q_i)}, sells S = {(pi_j, q_j)},
        tick grid P, tie-break cascade C, rationing rule R,
        reference price p_ref
OUTPUT: clearing price p*, fills F, residual E

1  candidates <- distinct limit prices in B and S, snapped to P
2  for each p in candidates:
3      D(p) <- sum of q_i over buys with pi_i >= p
4      S(p) <- sum of q_j over sells with pi_j <= p
5      V(p) <- min(D(p), S(p))
6  Vmax  <- max over p of V(p)
7  M     <- { p : V(p) = Vmax }              // maximising set
8  if Vmax = 0: return (no clearing, all orders roll over)
9  M     <- argmin over p in M of |D(p) - S(p)|      // cascade 1
10 if |M| > 1: M <- side of residual imbalance       // cascade 2
11 if |M| > 1: M <- argmin over p in M of |p - p_ref| // cascade 3
12 p*    <- unique element of M
13 fill all buys with pi_i > p* completely            // inframarginal
14 fill all sells with pi_j < p* completely
15 marginal <- orders with pi = p* on the long side
16 allocate Vmax - (filled inframarginal) among marginal using R
17 E     <- unmatched quantity on the long side
18 apply time-in-force to E: roll over, or cancel
19 return (p*, F, E)
\end{lstlisting}

Steps 1--7 are the line search; the unimodality of $V$ noted in
\Cref{sec:clearing} guarantees that the maximising set is an interval on
the grid. Steps 9--11 are the documented tie-break cascade, and step 16
is the only place where the rationing rule enters, which is why the
rules of \Cref{tab:rationing} can be swapped without touching the price
computation.

\chapter{Engine Configuration Parameters}
\label{app:config}

\begin{table}[htbp]
\centering
\caption[Engine configuration parameters]{Configuration parameters of
the two engines. The baseline column gives the values used in the
headline comparison of \Cref{ch:results}.}
\label{tab:config}
\small
\begin{tabularx}{\textwidth}{@{}p{3.6cm} L p{2.8cm}@{}}
\toprule
\textbf{Parameter} & \textbf{Meaning} & \textbf{Baseline}\\
\midrule
\texttt{engine}            & CDA or FBA & both, paired\\
\texttt{tau\_ns}           & Batch interval in nanoseconds & 100\,ms\\
\texttt{boundary}          & Fixed, randomised, or randomised extension
& fixed\\
\texttt{boundary\_seed}    & Seed for randomised closing & n/a\\
\texttt{rationing}         & Price-time, pro rata, size priority,
random, hybrid & pro rata\\
\texttt{hybrid\_kappa}     & Share allocated pro rata under the hybrid
rule & n/a\\
\texttt{tie\_break}        & Cascade order for the maximising set &
imbalance, side, reference\\
\texttt{residual\_policy}  & Roll over or cancel unmatched remainder &
roll over\\
\texttt{self\_match}       & Self-trade prevention policy & cancel
newest\\
\texttt{warmup\_ns}        & Lead-in excluded from metrics & 60\,s\\
\texttt{log\_level}        & Granularity of the structured log & full\\
\bottomrule
\end{tabularx}
\end{table}

\chapter{Reproducibility}
\label{app:repro}

Every run writes a log header containing the engine identity, the full
configuration of \Cref{tab:config}, the input file hash, and the code
commit hash. Given these, a run is byte-reproducible. The analysis
pipeline reads only the logs and is shared by both engines, so no metric
is computed by engine-specific code.

\ph[repository URL, dataset DOI and commit hash to be inserted before
submission]
